% ==== P4 draft O3 — §15.1 대체 블록 ====
\subsection{왜 흑연에는 없고 LCO 에는 있는가 --- MIT 의 물리}\label{sec:lco-why}
삽입 반응의 전체 엔트로피 변화는 ``리튬 자리 배열''(config)$\cdot$``격자 진동''(vib)$\cdot$``전자 준위 점유''
(electronic) 세 자유도의 합이며, 흑연에서는 셋째 몫이 작다 --- 흑연$\cdot$$\mathrm{LiC_6}$ 모두 전도성이 좋아
충방전 동안 $g(E_F)$ 가 크게 바뀌지 않아 전자 분포의 엔트로피 \emph{변화}가 거의 없다. LCO 는 다르다:
$x{=}1$($\mathrm{LiCoO_2}$)은 Co$^{3+}$($t_{2g}^6$ low-spin, 닫힌 껍질)이라 $t_{2g}$ 띠가 가득 차 그 위 빈
띠와 갭으로 갈린 \emph{밴드 절연체}(band insulator)이며($g(E_F)\!\approx\!0$), 단일입자 띠 그림만으로 이미
절연이다. 탈리튬화로 $x$ 가 $\sim$0.94 아래로 내려가면 Co$^{4+}$ 정공이 생겨 $t_{2g}$ 띠에 전도 전자가 열려
\emph{금속}이 된다($g(E_F)\!\to\!$유한)\cite{menetrier1999,motohashi2009}. 다만 강체 띠 채움 논리라면 정공이
열리는 순간 곧바로 금속이어야 하나, 실측 MIT 는 $x\approx0.75$--$0.94$ 의 2상역에서 일어나는 \emph{1차 전이}
이고(\S\ref{sec:lco-hys} 의 T1), 그 미시 해소는 전자 상관(correlation) 물리의 몫이다. 어느 기작이든 전자
분포는 이 2상역을 가로질러서만 켜지므로 전자 엔트로피 \emph{변화}도 그 구간에 국소한다.

\begin{bgbox}
\textbf{금속--절연체 전이와 Mott 물리 ---} 절연체는 단일입자 띠 그림만으로 설명되는가에 따라 두 부류로
갈린다. \emph{밴드 절연체}(band insulator)는 띠가 가득 차고 그 위 빈 띠와 갭으로 갈려, 전자를 하나하나
독립으로 본 띠 그림 안에서 이미 전도가 막힌다. \emph{Mott 절연체}(Mott insulator)는 띠가 반쯤 차 그
그림으로는 금속이어야 하나, 한 자리에 두 전자가 겹칠 때의 Coulomb 반발 $U$ 가 이웃으로 건너뛰는
\emph{이동적분}(hopping integral) $t$ 를 눌러 전자가 제자리에 묶이는 \emph{상관}(correlation) 효과의
절연이다\cite{mott1968,imada1998}. 금속--절연체 전이(MIT)는 채움(도핑)$\cdot$띠폭(압력)$\cdot$상관 셋의
경쟁으로 일어난다\cite{imada1998}. $\mathrm{Li}_x\mathrm{CoO_2}$ 에서 $x{=}1$ 은 Co$^{3+}$ 의 $t_{2g}^6$
닫힌 껍질이 $t_{2g}$ 띠를 가득 채운 \emph{밴드 절연체}이지, 상관이 전자를 자리에 묶는 Mott 절연체가
아니다\cite{menetrier1999}. 강체 띠 채움으로 읽으면 탈리튬화가 정공을 여는 순간 금속화를 예상시키나,
관측된 MIT 는 $x\approx0.75$--$0.94$ 2상역의 \emph{1차 전이}다\cite{menetrier1999,reimers1992,motohashi2009}.
Marianetti$\cdot$Kotliar$\cdot$Ceder 는 정공이 Li 빈자리에 속박된 \emph{불순물 준위}(impurity level)를 만들고
그 불순물 띠가 임계 정공 농도에서 비국재화하는 \emph{불순물 Mott 전이}(impurity Mott transition)로 이 1차
MIT 와 2상 공존을 설명했다(대형 초격자 계산에서 $x{>}0.95$ 절연$\cdot$$x{<}0.75$ 금속)\cite{marianetti2004}
--- 상관 물리가 host $t_{2g}$ 띠가 아니라 불순물 준위에서 작동한다는 그림이다. forward 모델은 이 미시 기작
전체를 조성축 게이트 $g(E_F,x)$ 하나로 현상학화하여, 게이트 중심을 전이 조성(실측 2상역의 중앙)으로, 폭을
발현 조성의 결함 분산으로 잡는다(게이트 형태의 정당화는 \S\ref{sec:lco-gate}). 곧 미시 기작은 전이가 왜
급격한 1차인지를 답하고, 게이트는 상태밀도 $g(E_F,x)$ 가 조성에 따라 어떻게 켜지는지를 담는다.
\end{bgbox}

곧 전자항이 이 전이 구간에서만 켜지는 게이트는 MIT 의 직접 결과이며, \S\ref{sec:lco-Se} 가 그 켜짐을
Fermi--Dirac 분포에서 Sommerfeld 전자 엔트로피로 닫는다.
% ==== 블록 끝 ====
%
% NOTE — Read Coverage / 보존 확인 / 자체검수
%
% [Read Coverage] 전문 정독:
%   - _sections/ch1_sec15_lcoelec.tex 전문(§15.0 도입~§15.5, 특히 §15.1 원문 L12–20 및 §15.4
%     게이트 정당화 sec:lco-gate — bgbox=왜(기작)/§15.4=게이트 형태 정당화 역할 분담 확인).
%   - _sections/ch1_sec13_lcohys.tex L120–164(T1 슬롯 분리 sec:lco-hys-mit — L136–137 이 미시
%     Mott 기작·전자 자유도 배경을 sec:lco-why 로 위임함을 확인, 본 bgbox 가 그 위임의 이행).
%   - results/V1020_STYLE_RUBRIC.md 전문(A1·A2 D1/D2 오독가드·A4·C1 병기·C4 두문자어·D2' 자족
%     bgbox·D3' 배경다리·E1 등재키만).
%   - MIT 두문자어 첫 확장 위치 확인: ch1_sec11_lcointro.tex L42 "insulator-to-metal transition, MIT"
%     — §15 는 첫 출현이 아니므로 본문은 bare MIT 유지(C1/C4 준수). bgbox 는 자족 요건상 한글
%     전체어 "금속--절연체 전이(MIT)" 로 자체 표기(D2').
%
% [보존 확인] 기존 §15.1 문장·인용 전량 유지:
%   (1) 세 자유도(config·vib·electronic) 합 서술 — 원문 그대로.
%   (2) 흑연 셋째 몫 작음(흑연·LiC6 전도성·g(E_F) 불변·엔트로피 변화 거의 없음) — 원문 그대로.
%   (3) x=1 LiCoO2 = Co3+ t2g^6 low-spin 닫힌 껍질, g(E_F)≈0 — 유지(+"전기 절연체"를 물리 부류
%       "밴드 절연체"로 정밀화. 절연 사실·g(E_F)≈0 보존).
%   (4) 탈리튬화 x<~0.94 → Co4+ 정공 → t2g 전도 → 금속 g(E_F)→유한 \cite{menetrier1999,motohashi2009}
%       — 문장·인용 원문 그대로 유지.
%   (5) MIT x≈0.75–0.94 2상역, \S\ref{sec:lco-hys} 의 T1 상호참조 — 유지.
%   (6) 그 구간에서만 전자 분포 켜짐 → 전자 엔트로피 변화 국소 — 유지("어느 기작이든" 로 기작-무관 강화).
%   (7) "게이트는 MIT 의 직접 결과다" — 마무리 다리에 보존(+§15.2 sec:lco-Se 로 전달).
%   인용 유실 0: menetrier1999·motohashi2009 본문 유지, bgbox 에 reimers1992·mott1968·imada1998·
%   marianetti2004 추가(브리프 6키 내). 신규 식·라벨 0.
%
% [자체검수 — 물리 하드 가드 3항]
%   (G1) LiCoO2(x=1) Mott 절연체 오기 금지 → 준수. 본문·bgbox 모두 "밴드 절연체(t2g^6 가득 찬 띠)"로
%        명시, bgbox 에 "밴드 절연체이지 Mott 절연체가 아니다"(A2 허용 오독가드 대조형)로 못박음.
%   (G2) Marianetti 과대해석 금지 → 준수. "설명했다"/"~작동한다는 그림이다" 수준(확정된 유일 기작 표현
%        없음). 계산 수치(x>0.95 절연·x<0.75 금속)는 "대형 초격자 계산에서"로 계산 소관 명시(실측과 구분).
%   (G3) Marianetti Mott 물리는 host 띠 아닌 불순물 준위 소관 → 준수. "상관 물리가 host t2g 띠가 아니라
%        불순물 준위에서 작동한다"로 명시.
%   (기타) §13 T1 슬롯분리 무모순: 본 소절은 전자 자유도 배경·게이트만 다루고 히스 gap(Ω1)·이중계산
%        경계는 건드리지 않음. §15.4 역할중복 회피: bgbox 는 기작(왜)만, 게이트 형태·3초기값 정당화는
%        \S\ref{sec:lco-gate} 로 위임(중심·폭 대응만 1문장 다리로 접속).
%   (자족성) bgbox 는 본문 지시대명사 무참조(내부 지시대명사만) — 부록 이설 시 성립. 본문은 bgbox 제거 시
%        본문→마무리 다리로 그대로 읽힘("미시 해소는 상관 물리의 몫" 자기완결 서술).
%
% 자산: [V20-096] [V20-097] [V20-098] [V20-099]
